\textbf{Phần 1. Trắc nghiệm lựa chọn}
\setcounterpageref{ex}{0}

\begin{ex} %[4D8N39]
Kết quả của phép chia $84\,628 : 2$ là
\choice
{\True $42\,314$}
{$42\,313$}
{$42\,304$}
{$41\,314$}
\loigiai{
$84\,628 : 2 = 42\,314$.\\
\textbf{Đáp án của câu hỏi là A}
}
\end{ex}

\begin{ex} %[4D8N39]
Số dư của phép chia $245\,678 : 5$ là
\choice
{$1$}
{$2$}
{\True $3$}
{$4$}
\loigiai{
$8 : 5 = 1$ dư 3.\\
\textbf{Đáp án của câu hỏi là C}
}
\end{ex}

\begin{ex} %[4D8H39]
Kết quả của phép chia $360\,000 : 9$ là
\choice
{$4000$}
{\True $40\,000$}
{$400\,000$}
{$30\,000$}
\loigiai{
$360\,000 : 9 = 40\,000$.\\
\textbf{Đáp án của câu hỏi là B}
}
\end{ex}

\begin{ex} %[4D8H39]
Tìm $x$ biết: $x \times 4 = 184\,000$.
\choice
{$45\,000$}
{$47\,000$}
{$48\,000$}
{\True $46\,000$}
\loigiai{
$x = 184\,000 : 4 = 46\,000$.\\
\textbf{Đáp án của câu hỏi là D}
}
\end{ex}

\begin{ex} %[4D8H39]
Một hình vuông có chu vi $96\text{ cm}$. Độ dài cạnh hình vuông là
\choice
{$22\text{ cm}$}
{\True $24\text{ cm}$}
{$26\text{ cm}$}
{$48\text{ cm}$}
\loigiai{
$96 : 4 = 24\text{ cm}$.\\
\textbf{Đáp án của câu hỏi là B}
}
\end{ex}

\begin{ex} %[4D8H39]
Kết quả của phép chia $125\,000 : 5$ là
\choice
{\True $25\,000$}
{$20\,000$}
{$30\,000$}
{$250\,000$}
\loigiai{
$125\,000 : 5 = 25\,000$.\\
\textbf{Đáp án của câu hỏi là A}
}
\end{ex}

\begin{ex} %[4D8V39]
Người ta đóng $12\,450\text{ kg}$ gạo vào các bao loại $5\text{ kg}$. Số bao gạo đóng được là
\choice
{$2480$ bao}
{$2500$ bao}
{$2450$ bao}
{\True $2490$ bao}
\loigiai{
$12\,450 : 5 = 2490$ bao.\\
\textbf{Đáp án của câu hỏi là D}
}
\end{ex}

\begin{ex} %[4D8V39]
Có $350$ học sinh đi tham quan bằng xe ô tô loại 45 chỗ ngồi (không tính ghế lái xe). Cần ít nhất bao nhiêu xe ô tô như thế để chở hết số học sinh?
\choice
{$7$ xe}
{$9$ xe}
{\True $8$ xe}
{$10$ xe}
\loigiai{
$350 : 45 = 7$ dư 35 em. Cần ít nhất $7 + 1 = 8$ xe.\\
\textbf{Đáp án của câu hỏi là C}
}
\end{ex}

\begin{ex} %[4D8V39]
Một xí nghiệp sản xuất được $48\,600$ hộp bánh trong 6 ngày. Hỏi trung bình mỗi ngày xí nghiệp sản xuất được bao nhiêu hộp bánh?
\choice
{\True $8100$ hộp}
{$8000$ hộp}
{$8200$ hộp}
{$8300$ hộp}
\loigiai{
$48\,600 : 6 = 8100$ hộp.\\
\textbf{Đáp án của câu hỏi là A}
}
\end{ex}

\begin{ex} %[4D8V39]
Tính giá trị biểu thức: $450\,000 - 360\,000 : 4$.
\choice
{$340\,000$}
{$380\,000$}
{$22\,500$}
{\True $360\,000$}
\loigiai{
$450\,000 - 90\,000 = 360\,000$.\\
\textbf{Đáp án của câu hỏi là D}
}
\end{ex}

\begin{ex} %[4D8V39]
Tìm số bị chia trong phép chia cho 7 có thương là $12\,340$ và số dư lớn nhất có thể.
\choice
{$86\,380$}
{$86\,385$}
{\True $86\,386$}
{$86\,387$}
\loigiai{
Số dư lớn nhất là 6.\\
Số bị chia: $12\,340 \times 7 + 6 = 86\,380 + 6 = 86\,386$.\\
\textbf{Đáp án của câu hỏi là C}
}
\end{ex}

\begin{ex} %[4D8C39]
Khi chia một số cho 6 được thương là $15\,200$ và không còn dư. Nếu chia số đó cho 8 thì được thương là bao nhiêu?
\choice
{$11\,300$}
{\True $11\,400$}
{$11\,500$}
{$12\,000$}
\loigiai{
Số đó là: $15\,200 \times 6 = 91\,200$.\\
Chia cho 8: $91\,200 : 8 = 11\,400$.\\
\textbf{Đáp án của câu hỏi là B}
}
\end{ex}

\textbf{Phần 2. Trắc nghiệm đúng sai}
\setcounterpageref{ex}{0}

\begin{ex} %[4D8H39]
Xét tính đúng sai của các phép chia:
\choiceTF
{\True $150\,000 : 3 = 50\,000$.}
{$45\,678 : 9 = 5075$ dư 5.}
{\True $45\,678 : 9 = 5075$ dư 3.}
{\True Trong phép chia có dư, số dư luôn nhỏ hơn số chia.}
\loigiai{
\begin{itemize}
    \item $150\,000 : 3 = 50\,000$. Đúng.
    \item $45\,678 : 9 = 5075$ dư 3 ($5075 \times 9 + 3 = 45678$). Ý 2 Sai, Ý 3 Đúng.
    \item Số dư $<$ số chia: Đúng.
\end{itemize}
\begin{itemize}
    \item a) \textbf{Đúng}.
    \item b) \textbf{Sai}.
    \item c) \textbf{Đúng}.
    \item d) \textbf{Đúng}.
\end{itemize}
}
\end{ex}

\begin{ex} %[4D8V39]
Có $120$ lít mật ong đóng đều vào các can $5$ lít. Mỗi can bán với giá $450\,000$ đồng. Xét tính đúng sai:
\choiceTF
{\True Đóng được tất cả $24$ can mật ong.}
{\True Mỗi lít mật ong có giá là $90\,000$ đồng.}
{Tổng số tiền bán mật ong là $10\,000\,000$ đồng.}
{\True Tổng số tiền bán mật ong là $10\,800\,000$ đồng.}
\loigiai{
\begin{itemize}
    \item $120 : 5 = 24$ can. Đúng.
    \item $450\,000 : 5 = 90\,000$ đồng/lít. Đúng.
    \item Tổng tiền: $24 \times 450\,000 = 10\,800\,000$ đồng $\neq 10\,000\,000$ đồng. Ý 3 Sai, Ý 4 Đúng.
\end{itemize}
\begin{itemize}
    \item a) \textbf{Đúng}.
    \item b) \textbf{Đúng}.
    \item c) \textbf{Sai}.
    \item d) \textbf{Đúng}.
\end{itemize}
}
\end{ex}

\textbf{Phần 3. Trả lời ngắn}
\setcounterpageref{ex}{0}

\begin{ex} %[4D8H39]
Tính thương: $240\,000 : 6$.
\shortans{40000}
\loigiai{
\textbf{Đáp số:} $40000$\\
$40\,000$.
}
\end{ex}

\begin{ex} %[4D8H39]
Tính thương: $500\,000 : 5$.
\shortans{100000}
\loigiai{
\textbf{Đáp số:} $100000$\\
$100\,000$.
}
\end{ex}

\begin{ex} %[4D8V39]
Tìm $x$ biết: $x \times 8 = 320\,000$.
\shortans{40000}
\loigiai{
\textbf{Đáp số:} $40000$\\
$x = 320\,000 : 8 = 40\,000$.
}
\end{ex}

\begin{ex} %[4D8V39]
Một cuộn dây dài $450\text{ m}$, người ta cắt thành các đoạn dài $9\text{ m}$. Hỏi cắt được bao nhiêu đoạn dây?
\shortans{50}
\loigiai{
\textbf{Đáp số:} $50$\\
$450 : 9 = 50$ đoạn.
}
\end{ex}

\begin{ex} %[4D8V39]
Có $155\text{ kg}$ gạo đóng vào các túi $5\text{ kg}$. Hỏi đóng được bao nhiêu túi?
\shortans{31}
\loigiai{
\textbf{Đáp số:} $31$\\
$155 : 5 = 31$ túi.
}
\end{ex}

\begin{ex} %[4D8C39]
Một số chia cho 5 dư 3. Hỏi số đó cộng thêm ít nhất bao nhiêu đơn vị để được số chia hết cho 5?
\shortans{2}
\loigiai{
\textbf{Đáp số:} $2$\\
$5 - 3 = 2$ đơn vị.
}
\end{ex}
